\chapter{Continuations}
\index{continuations}\index{call/cc}
\label{ch:continuations}

Continuations are Scheme's most distinctive feature. A continuation represents ``the rest of the computation''---everything that would happen after a particular point in the program. Scheme lets you capture this as a first-class value and invoke it later, enabling non-local exits, coroutines, backtracking, and custom control structures.

This chapter starts with simple escape continuations (early returns), then covers full continuations, and finishes with \texttt{dynamic-wind} for cleanup guarantees.

\section{What Is a Continuation?}

At every point during execution, there is a ``continuation''---the rest of the work to be done. Consider:

\begin{lstlisting}
(* 3 (+ 1 2))
\end{lstlisting}

When Scheme evaluates \texttt{(+ 1 2)}, the continuation is ``take the result and multiply it by 3.'' Normally, continuations are implicit---you never see them. Scheme's \texttt{call/cc} makes them explicit.

\section{Escape Continuations with \texttt{call/ec}}
\index{call/ec}\index{escape continuation}

The most common use of continuations is an \emph{early return}---exiting a computation before it would normally finish. \texttt{call/ec} (call with escape continuation) provides this:

\begin{lstlisting}[style=repl]
> (call/ec (lambda (return)
             (return 42)
             (display "never reached")))
42
\end{lstlisting}

When \texttt{return} is invoked, execution immediately jumps out of the \texttt{call/ec} form with the given value. Code after the call is skipped.

\subsection{Early Return from a Loop}

A common pattern: searching a list and returning the first match:

\begin{lstlisting}[style=repl]
> (define (find-negative lst)
    (call/ec (lambda (return)
      (for-each (lambda (x)
                  (when (negative? x) (return x)))
                lst)
      #f)))
> (find-negative '(3 1 -4 1 5))
-4
> (find-negative '(1 2 3))
#f
\end{lstlisting}

Without \texttt{call/ec}, you would need recursion or a flag variable. The escape continuation gives you the equivalent of Python's \texttt{return} inside a loop.

\subsection{Breaking Out of Nested Computation}

\begin{lstlisting}[style=repl]
> (define (product lst)
    (call/ec (lambda (bail)
      (let loop ((rest lst) (acc 1))
        (cond
          ((null? rest) acc)
          ((zero? (car rest)) (bail 0))
          (else (loop (cdr rest) (* acc (car rest)))))))))
> (product '(1 2 3 4 5))
120
> (product '(1 2 0 4 5))
0
\end{lstlisting}

If we see a zero, we bail immediately---no point multiplying the remaining elements.

\begin{note}[Coming from Python]
\texttt{call/ec} is like wrapping code in a function and using \texttt{return}. Scheme does not have a \texttt{return} statement, so \texttt{call/ec} fills that role. The important difference: it works across function boundaries (the escape continuation can be invoked from inside a lambda passed to \texttt{for-each}).
\end{note}

\section{Full Continuations with \texttt{call/cc}}
\index{call/cc}

\texttt{call/cc} (short for \texttt{call-with-current-continuation}) captures the \emph{full} continuation---not just an escape, but a snapshot of the entire execution state that can be invoked multiple times, even after the original computation has finished.

\begin{lstlisting}[style=repl]
> (call/cc (lambda (k) (k 42)))
42
> (+ 1 (call/cc (lambda (k) (+ 2 (k 3)))))
4
\end{lstlisting}

In the second example, the continuation of \texttt{call/cc} is ``add 1 to the result.'' Calling \texttt{(k 3)} jumps to that point with 3, so the answer is 4. The \texttt{(+ 2 ...)} is abandoned.

\subsection{Saving and Reusing a Continuation}

The real power of \texttt{call/cc} is that you can save the continuation and invoke it later:

\begin{lstlisting}[style=repl]
> (define saved #f)
> (define (capture)
    (call/cc (lambda (k)
      (set! saved k)
      "first time")))
> (capture)
"first time"
> (saved "second time")
"second time"
> (saved "third time")
"third time"
\end{lstlisting}

Each call to \texttt{saved} re-enters the computation at the point where \texttt{call/cc} was called, as if it had returned the given value.

\subsection{A Counter with Continuations}

\begin{lstlisting}[style=repl]
> (define resume #f)
> (define (start)
    (let ((result (call/cc (lambda (k)
                    (set! resume k)
                    0))))
      (display result)
      (display " ")
      (when (< result 3)
        (resume (+ result 1)))))
> (start)
0 1 2 3
\end{lstlisting}

Each call to \texttt{resume} re-enters the \texttt{let}, with \texttt{result} bound to the new value. This is how coroutines and generators can be implemented.

\section{\texttt{call/ec} vs.\ \texttt{call/cc}}

\begin{center}
\begin{tabular}{lp{5cm}l}
\textbf{Feature} & \textbf{\texttt{call/ec}} & \textbf{\texttt{call/cc}} \\
\hline
Can be invoked after proc returns & No & Yes \\
Can be invoked multiple times & No & Yes \\
Stack snapshot required & No & Yes \\
Performance & Fast (no allocation) & Slow (copies stack) \\
Use case & Early return, loop break & Coroutines, backtracking \\
\end{tabular}
\end{center}

\begin{warning}[Prefer \texttt{call/ec}]
Use \texttt{call/ec} unless you genuinely need to reinvoke the continuation after the procedure returns. Full \texttt{call/cc} copies the entire call stack, which is O(stack depth). For early exits and error handling, \texttt{call/ec} is dramatically cheaper.
\end{warning}

\section{\texttt{dynamic-wind}: Guaranteed Cleanup}
\index{dynamic-wind}

\texttt{dynamic-wind} ensures that setup and cleanup code run whenever control enters or leaves a region---even when continuations jump in or out:

\begin{lstlisting}
(dynamic-wind
  before-thunk    ; called when entering
  body-thunk      ; the main computation
  after-thunk)    ; called when leaving
\end{lstlisting}

\begin{lstlisting}[style=repl]
> (let ((log '()))
    (dynamic-wind
      (lambda () (set! log (cons 'enter log)))
      (lambda () (set! log (cons 'body log)))
      (lambda () (set! log (cons 'leave log))))
    (reverse log))
(enter body leave)
\end{lstlisting}

\subsection{Resource Cleanup}

The most common use is ensuring resources are released:

\begin{lstlisting}
(define (with-open-file filename proc)
  (let ((port (open-input-file filename)))
    (dynamic-wind
      (lambda () #f)
      (lambda () (proc port))
      (lambda () (close-input-port port)))))
\end{lstlisting}

The port is closed whether the body returns normally, raises an exception, or a continuation exits.

\subsection{Re-entry Behavior}

If a continuation captured inside the body is invoked from outside, the ``before'' thunk runs again. If a continuation from outside is invoked from inside, the ``after'' thunk runs:

\begin{lstlisting}[style=repl]
> (let ((log '()) (k #f))
    (dynamic-wind
      (lambda () (set! log (cons 'in log)))
      (lambda ()
        (call/cc (lambda (c) (set! k c)))
        (set! log (cons 'body log)))
      (lambda () (set! log (cons 'out log))))
    (when (< (length log) 6) (k #f))
    (reverse log))
(in body out in body out)
\end{lstlisting}

The continuation re-enters, triggering the ``in'' thunk again, then the body runs again, then ``out.''

\section{Practical Uses of Continuations}

\subsection{Implementing Exceptions}

Before \texttt{guard} existed, exceptions were implemented with continuations:

\begin{lstlisting}
(define (try thunk handler)
  (call/cc (lambda (exit)
    (with-exception-handler
      (lambda (e) (exit (handler e)))
      thunk))))
\end{lstlisting}

\subsection{Coroutines}

A simple coroutine system using continuations:

\begin{lstlisting}
(define (make-coroutine proc)
  (let ((resume #f))
    (define (yield value)
      (call/cc (lambda (k)
        (set! resume k)
        value)))
    (lambda ()
      (if resume
          (resume #f)
          (proc yield)))))
\end{lstlisting}

\subsection{Non-Deterministic Choice (Backtracking)}

\begin{lstlisting}
(define fail #f)

(define (choose . alternatives)
  (call/cc (lambda (return)
    (let loop ((alts alternatives))
      (if (null? alts)
          (fail)
          (let ((old-fail fail))
            (call/cc (lambda (k)
              (set! fail (lambda ()
                (set! fail old-fail)
                (k #f)))
              (return (car alts))))
            (loop (cdr alts))))))))
\end{lstlisting}

This creates a backtracking mechanism: \texttt{choose} returns each alternative in turn, backtracking when \texttt{fail} is called.

\section{When to Use Continuations}

Continuations are powerful but should be used sparingly:

\begin{itemize}
    \item \textbf{Use \texttt{call/ec}} for early returns, loop breaks, and search. This is common and cheap.
    \item \textbf{Use \texttt{guard}} for exception handling. It is built on continuations but provides a cleaner interface.
    \item \textbf{Use \texttt{call/cc}} only for advanced patterns: coroutines, generators, backtracking, or implementing new control structures. These are rare in application code.
    \item \textbf{Always pair with \texttt{dynamic-wind}} when resources need cleanup across continuation invocations.
\end{itemize}

Most programs never need full \texttt{call/cc}. But knowing it exists---and that all control flow in Scheme is ultimately built on continuations---gives you a deep understanding of the language's power.

\section*{Exercises}

\begin{exercise}[Find-First with call/ec]
Use \code{call/ec} to write a \code{find-first} procedure that takes a predicate and a list, and returns the first element satisfying the predicate, or \code{\#f} if none does. Do not use the built-in \code{find}---use \code{for-each} with an escape continuation instead.

\begin{lstlisting}[style=repl]
> (find-first even? '(1 3 4 7 8))
4
> (find-first negative? '(1 2 3))
#f
> (find-first string? '(1 "hello" 3))
"hello"
\end{lstlisting}
\end{exercise}

\begin{exercise}[Early Exit from Nested Loops]
Use \code{call/ec} to write a \code{find-pair-summing-to} procedure that takes a target number and a list, and returns the first pair of elements whose sum equals the target (as a list of two elements), or \code{\#f} if no such pair exists. This is like Python's nested loop with an early \code{return}.

\begin{lstlisting}[style=repl]
> (find-pair-summing-to 10 '(1 3 7 4 6))
(3 7)
> (find-pair-summing-to 100 '(1 2 3))
#f
\end{lstlisting}

Hint: nest two \code{for-each} calls inside a single \code{call/ec}.
\end{exercise}

\begin{exercise}[Resource Cleanup with dynamic-wind]
Write a \code{with-resource} procedure that takes three arguments: an \emph{acquire} thunk (returns a resource), a \emph{body} procedure (receives the resource and uses it), and a \emph{release} procedure (cleans up the resource). Use \code{dynamic-wind} to guarantee the release procedure runs even if an error is raised.

\begin{lstlisting}[style=repl]
> (with-resource
    (lambda () (display "acquiring\n") 'my-resource)
    (lambda (r) (display "using ") (display r) (newline) 42)
    (lambda (r) (display "releasing ") (display r) (newline)))
acquiring
using my-resource
releasing my-resource
42
\end{lstlisting}

Verify that the release runs even when the body raises an error.
\end{exercise}

\begin{exercise}[Challenge: Simple Generator]
Implement a generator using \code{call/cc}. Write a \code{make-generator} procedure that takes a procedure. The procedure receives a \code{yield} function; each call to \code{yield} suspends the generator and returns the yielded value. Calling the generator again resumes from where it left off. When the procedure returns, the generator returns \code{'done} on all subsequent calls.

\begin{lstlisting}[style=repl]
> (define gen (make-generator
    (lambda (yield)
      (yield 1)
      (yield 2)
      (yield 3))))
> (gen)
1
> (gen)
2
> (gen)
3
> (gen)
done
\end{lstlisting}

This mirrors Python's \texttt{yield} statement but is built entirely from \code{call/cc}.
\end{exercise}

You now understand both escape continuations (practical, efficient) and full continuations (powerful, rare). Together with \texttt{dynamic-wind}, these tools give Scheme its unique ability to express any control flow pattern. This completes Part~III. In Part~IV, we explore features beyond the R7RS standard---the SRFI ecosystem, the foreign function interface, and concurrency.

\chapterend
